\let\negmedspace\undefined
\let\negthickspace\undefined
\documentclass[journal]{IEEEtran}
\usepackage[a5paper, margin=10mm, onecolumn]{geometry}
\usepackage{tfrupee}
\setlength{\headheight}{1cm}
\setlength{\headsep}{0mm}
\usepackage{gvv-book}
\usepackage{gvv}
\usepackage{cite}
\usepackage{amsmath,amssymb,amsfonts,amsthm}
\usepackage{algorithmic}
\usepackage{graphicx}
\usepackage{textcomp}
\usepackage{xcolor}
\usepackage{txfonts}
\usepackage{listings}
\usepackage{enumitem}
\usepackage{mathtools}
\usepackage{gensymb}
\usepackage{comment}
\usepackage[breaklinks=true]{hyperref}
\usepackage{multicol}

\begin{document}

\bibliographystyle{IEEEtran}
\vspace{2cm}

\title{5.5.25}
\author{AI25BTECH11008 - Chiruvella Harshith Sharan}
{\let\newpage\relax\maketitle}

\textbf{Question:} \\

If
\[
A = \myvec{
1 & 2 & -3\\
2 & 0 & -3\\
1 & 2 & 0
}
\]
then find \(A^{-1}\) and use it to solve the system
\[
x + 2y - 3z = 1,\qquad 2x - 3z = 2,\qquad x + 2y = 3.
\]

\vspace{0.3cm}
\textbf{Solution:}

\subsection*{Step 1: Matrix form}
Write the system as a matrix equation
\[
A\vec{x} = \vec{b}. \tag{1}
\]
Here
\[
\vec{x} = \myvec{x\\y\\z},\qquad \vec{b}=\myvec{1\\2\\3}. \tag{2}
\]

\subsection*{Step 2: Compute \(\det(A)\) and the adjugate}
We compute the determinant by expansion (first row):
\[
\det(A)=
1\cdot\det\myvec{0 & -3\\[2pt] 2 & 0}
-2\cdot\det\myvec{2 & -3\\[2pt] 1 & 0}
-3\cdot\det\myvec{2 & 0\\[2pt] 1 & 2}
=1\cdot 6 -2\cdot 3 -3\cdot 4 = -12.
\]

Next compute the matrix of cofactors (each entry is the signed minor). The cofactor matrix \(C\) is
\[
C = \myvec{
\;6 & -3 & \;4\\[4pt]
-6 & \;3 & \;0\\[4pt]
-6 & -3 & -4
}.
\]
The adjugate is the transpose of \(C\):
\[
\operatorname{adj}(A)=C^{\mathsf T}=
\myvec{
6 & -6 & -6\\[4pt]
-3 & 3 & -3\\[4pt]
4 & 0 & -4
}. \tag{3}
\]

\subsection*{Step 3: Inverse from adjugate and determinant}
Using \(A^{-1}=\dfrac{1}{\det(A)}\operatorname{adj}(A)\) we get
\[
A^{-1}=\frac{1}{-12}
\myvec{
6 & -6 & -6\\
-3 & 3 & -3\\
4 & 0 & -4
}
=
\myvec{
-\tfrac{1}{2} & \tfrac{1}{2} & \tfrac{1}{2}\\[4pt]
\tfrac{1}{4} & -\tfrac{1}{4} & \tfrac{1}{4}\\[4pt]
-\tfrac{1}{3} & 0 & \tfrac{1}{3}
}. \tag{4}
\]

\subsection*{Step 4: Solve \(\vec{x}=A^{-1}\vec{b}\)}
Multiply to obtain the solution:
\[
\vec{x}=A^{-1}\vec{b}
=
\myvec{
-\tfrac{1}{2} & \tfrac{1}{2} & \tfrac{1}{2}\\[4pt]
\tfrac{1}{4} & -\tfrac{1}{4} & \tfrac{1}{4}\\[4pt]
-\tfrac{1}{3} & 0 & \tfrac{1}{3}
}
\myvec{1\\2\\3}
=
\myvec{2\\[4pt]\tfrac{1}{2}\\[4pt]\tfrac{2}{3}}.
\]

\subsection*{Final Answer:}
\[
\boxed{x=2,\quad y=\tfrac{1}{2},\quad z=\tfrac{2}{3}}
\]

Below is a schematic TikZ illustration projecting the first two components of the vectors (a projection to the $xy$--plane for visualization). The algebraic relation is exact in $\mathbb{R}^3$; the drawing is only a 2D aid.
\begin{figure}[h!]
\centering \begin{tikzpicture}[scale=1.0]
% axes 
\draw[->] (-1,0) -- (5,0) node[right] {$x$};
\draw[->] (0,-4) -- (0,4) node[above] {$y$};
% column vectors (projected to xy-plane using first two components) 
\draw[->, very thick, blue] (0,0) -- (1,2) node[midway, above left] {$\mathbf{a}_1=(1,2,1)$}; \draw[->, very thick, green!60!black] (0,0) -- (2,0) node[midway, below] {$\mathbf{a}_2=(2,0,2)$}; \draw[->, very thick, red] (0,0) -- (-3,-3) node[midway, below left] {$\mathbf{a}_3=(-3,-3,0)$}; 
% scaled vectors (showing coefficients) 
\draw[->, dashed, blue] (0,0) -- (2,4) node[near end, right] {$2\mathbf{a}_1$}; 
\draw[->, dashed, green!60!black] (2,0) -- (3,0) node[midway, above] {$\tfrac{1}{2}\mathbf{a}_2$}; % this is from a2 tip for illustration
\draw[->, dashed, red] (0,0) -- (-2,-2) node[near end, left] {$\tfrac{2}{3}\mathbf{a}_3$};
% b (projected) 
\draw[->, thick, purple] (0,0) -- (1,2) node[midway, above right] {$\mathbf{b}=(1,2,3)$};
% legend (simple)
\node[anchor=west] at (3.2,3.2) {\small Linear combination: $2\mathbf{a}_1+\tfrac{1}{2}\mathbf{a}_2+\tfrac{2}{3}\mathbf{a}_3=\mathbf{b}$};
\node[anchor=west] at (3.2,2.6) {\small (projection shown for visualization only)};
\end{tikzpicture} 
\caption{Schematic (projection) showing column vectors of $A$ and the right-hand side $\mathbf{b}$.} 
\end{figure}

\end{document}



